\documentclass[11pt]{article}
\usepackage{icmlko}
\begin{document}

\title{같은 해 안에서도 버는가 --- 두 해 된 거부권이 잡지 못한 것을 잡는 자를, 알려진 날짜 구멍으로 검증했다}
\author{월드모델 실험실}
\date{노트 541 · 2026년 8월 4일}
\maketitle

\begin{abstract}
노트 446 거부권(\emph{판$\uparrow$ 인데 날짜 통제 판$\downarrow$ 이면 기각})은
``그 이득은 날짜 구멍''이라는 \textbf{읽기} 위에 서 있다. 노트 539 가 반례를
냈다 --- 시장팝업 \texttt{category} 는 연도를 안 실어 나르는데(크루스칼
$p = 0.168$) 그 무늬가 났다. \textbf{더 나은 자를 세우고 알려진 날짜 구멍으로
검증했다.}

\begin{center}
\begin{tabular}{lrr}
\hline
 & \textbf{같은 해 쌍} & 다른 해 쌍 \\
\hline
\textbf{TIMEAX}(알려진 날짜 구멍) & $\mathbf{-0.00189}$ & $+0.00125$ \\
 & $[-0.0025, -0.0013]$ & $[+0.0005, +0.0020]$ \\
\textbf{시장팝업 category} & $\mathbf{+0.00112}$ & $+0.00302$ \\
 & $[+0.0007, +0.0015]$ & $[+0.0026, +0.0035]$ \\
\hline
\end{tabular}
\end{center}

\textbf{새 자가 둘을 정확히 가른다} --- 날짜 구멍은 같은 해에서 \emph{잃고},
진짜 정보는 같은 해에서도 \emph{번다}. 연도는 예보 시점에 아는 값이므로
자로 쓸 수 있다.
\end{abstract}

\begin{figure}[t]
\centering
\includegraphics[width=\textwidth]{figs/f1.pdf}
\caption{\textbf{가} --- 두 팔. \textbf{나} --- 옛 자와 새 자.
\textbf{다} --- 내 대조가 처음엔 안 섰다.}
\end{figure}

\section{내 양성 대조가 처음엔 안 섰다}

먼저 현 챔피언 위에서 \texttt{TIMEAX} 를 켰다. 같은 해 $-0.0040$ $\cdot$
다른 해 $-0.0018$ 로 \textbf{둘 다 음수}가 나왔다 --- 날짜 구멍의 무늬가
아니다.

\begin{quote}
이유는 원장에 있었다. 노트 $524$ 가 확정한 대로 \textbf{\texttt{TRAINR}
(학습 쪽 최근성)이 채택된 뒤로 \texttt{TIMEAX} 는 어디서도 안 번다} ---
``최근성은 한 번만 팔린다''. \textbf{날짜 구멍의 무늬는 \texttt{TRAINR} 이
없던 노트 $447$ 설정에서만 난다.}
\end{quote}

거기서 다시 재니 같은 해 $\mathbf{-0.00189}$ $\cdot$ 다른 해
$\mathbf{+0.00125}$ --- \textbf{정확히 그 무늬}다.

\begin{quote}
\textbf{정직한 단서} --- 대조를 고친 것은 시장팝업 결과를 본 \emph{뒤}였다.
고친 근거는 노트 $524$ 의 기존 사실이지 원하는 답이 아니었지만,
\textbf{더 엄격했다면 대조 설정을 미리 못박았어야} 한다. 다음부터 양성
대조는 \emph{어느 챔피언 위에서} 도는지까지 사전 등록에 적는다.
\end{quote}

\section{옛 자는 둘을 못 가르고 새 자는 가른다}

\begin{center}
\begin{tabular}{lll}
\hline
 & 노트 446 거부권 & \textbf{같은 해 이득} \\
\hline
TIMEAX($447$ 설정) & 기각 (맞다) & $-0.0019$ $\to$ 기각 (맞다) \\
시장팝업 category & 기각 (\textbf{틀렸다}) & $\mathbf{+0.0011}$ $\to$ \textbf{통과} \\
\hline
\end{tabular}
\end{center}

거부권은 \emph{증상}(날짜 통제 판이 내려감)을 보고, 새 자는 \emph{원인}
(연도가 다른 쌍에서만 버는가)을 직접 본다. 시장팝업의 경우 날짜 통제 하락은
노트 539 가 보인 대로 \textbf{전부 이웃 도메인}이었고, 그 이웃 손해가
거부권을 오작동시켰다.

\section*{이 노트에서 규칙을 안 고친다}

사전 등록에 그렇게 적었고 지킨다. \textbf{사후에 규칙을 고쳐 통과시키면
그것이 점수 고르기다.} 규약 변경은 다음 노트에 따로 사전 등록한다 ---
바꿀 문장, 소급 적용 여부, 그리고 \emph{새 자가 틀릴 수 있는 방식}까지 적고
시작한다.

\section*{남은 것}

\begin{itemize}
\item \textbf{규약 변경 사전 등록}(노트 542) --- 거부권 446 을 ``판$\uparrow$
이면서 \emph{같은 해 이득}이 $0$ 이하이면 기각''으로 바꿀지.
\item \textbf{소급} --- 446 으로 기각된 과거 후보들을 새 자로 다시 볼지.
지금 원장에 몇 건인지부터 세야 한다.
\item \textbf{새 자의 한계} --- 유보가 $2025 \cdot 2026$ 두 해뿐이라 ``같은 해''가
쌍의 $50\!\sim\!72\%$ 다. 해가 늘면 이 자의 힘도 바뀐다.
\end{itemize}

\end{document}
